\providecommand{\Sp}{\mathrm{Sp}}
\providecommand{\Mp}{\mathrm{Mp}}
\providecommand{\Z}{\mathbb{Z}}
\providecommand{\Q}{\mathbb{Q}}
\providecommand{\R}{\mathbb{R}}
\providecommand{\C}{\mathbb{C}}
\providecommand{\CY}{\mathrm{CY}}
\providecommand{\SVOA}{\mathrm{SVOA}}
\providecommand{\VOA}{\mathrm{VOA}}
\providecommand{\NS}{\mathrm{NS}}
\providecommand{\ChirAlg}{\mathrm{ChirAlg}}
\providecommand{\Per}{\mathrm{Per}}
\providecommand{\Aut}{\mathrm{Aut}}
\providecommand{\id}{\mathrm{id}}
\providecommand{\Mod}{\mathrm{Mod}}
\providecommand{\kBKM}{\kappa_{\mathrm{BKM}}}
\providecommand{\Phisup}{\Phi^{\mathrm{sup}}}
\providecommand{\CYMp}{\CY_2^{\Mp}}
\providecommand{\CYint}{\CY_2^{\mathrm{int}}}
\providecommand{\CYhalf}{\CY_2^{\mathrm{half}}}
\providecommand{\cC}{\mathcal{C}}
\providecommand{\cK}{\mathcal{K}}
\providecommand{\cL}{\mathcal{L}}

\section{The metaplectic domain \(\CYMp\)}
\label{sec:CY2-Mp-domain}

\subsection*{Objects and cover data}

\paragraph{Definition.}
An object of \(\CYMp\) is a tuple
\[
  (\cC,L,\Gamma,\widetilde{\Gamma},\nu)
\]
with the following data.
\begin{itemize}
\item \(\cC\) is a smooth proper Kontsevich-Soibelman CY\(_2\)
  \((\infty,1)\)-category, with Serre functor the shift \([2]\).
\item \(L\subset K_0^{\mathrm{num}}(\cC)\otimes\R\) is a rank-four
  rational Mukai sublattice with symplectic Euler pairing.
\item \(\Gamma\subset\Sp_4(\Z)\) is the arithmetic monodromy group of
  the period map for \((\cC,L)\).
\item
  \(1\to\{\pm1\}\to\widetilde{\Gamma}\to\Gamma\to1\)
  is the metaplectic pullback of \(\Mp_4(\Z)\to\Sp_4(\Z)\).
\item \(\nu:\widetilde{\Gamma}\to\C^\times\) is a multiplier character
  specifying the automorphic line on which the intended Siegel or
  paramodular form lives.
\end{itemize}
The fibre over a CY\(_2\) category depends on the period monodromy
subgroup, the chosen arithmetic cover, and the multiplier.

\paragraph{Gerbe class.}
Let \(\cL_{\mathrm{per}}\) be the Hodge line over the period stack of
\((\cC,L)\).  The obstruction to descending the metaplectic lift to the
ordinary symplectic monodromy is
\[
  w_2^{\mathrm{per}}(\cC,L)\in
  H^2(\Per(\cC,L);\Z/2).
\]
The map \(\CYMp\to\CY_2\) is a \(B(\Z/2)\)-gerbe over the locus where
this obstruction is defined.  A splitting exists exactly after choosing
a trivialisation of \(w_2^{\mathrm{per}}\).  Without this datum the
notation \(\CY_2\times B\Mp_4\) describes only a split local model.

\subsection*{Morphisms}

\paragraph{Definition.}
A morphism
\[
  (\cC_1,L_1,\Gamma_1,\widetilde{\Gamma}_1,\nu_1)
  \longrightarrow
  (\cC_2,L_2,\Gamma_2,\widetilde{\Gamma}_2,\nu_2)
\]
is a CY\(_2\) Fourier-Mukai functor
\[
  F:\cC_1\to\cC_2,\qquad
  \cK_F\in\cC_1^{\mathrm{op}}\otimes\cC_2,
\]
such that \(F_\ast:L_1\to L_2\) is a Mukai isometry and the induced
period-map morphism sends the lifted monodromy and multiplier of the
source to those of the target.  Equivalently, the square
\[
\begin{array}{ccc}
\widetilde{\Gamma}_1 & \longrightarrow & \widetilde{\Gamma}_2\\
\downarrow && \downarrow\\
\Gamma_1 & \longrightarrow & \Gamma_2
\end{array}
\]
commutes and \(\nu_1=\nu_2\circ\widetilde{F}_\ast\).  Two morphisms are
equivalent when their Fourier-Mukai kernels are isomorphic and the
isomorphism has trivial Weil coboundary.

\paragraph{Composition.}
Composition is convolution of Fourier-Mukai kernels together with
multiplication of the corresponding Weil cocycles.  The associativity
constraint is the standard cocycle identity for the metaplectic central
extension.  Thus \(\CYMp\) is an \((\infty,1)\)-category after restricting
to kernels whose period action preserves the chosen lifted arithmetic
data.

\subsection*{Integer-weight subcategory}

\paragraph{Definition.}
\(\CYint\subset\CYMp\) is the full subcategory of objects for which the
multiplier is trivial on the central kernel of
\(\widetilde{\Gamma}\to\Gamma\).  Equivalently, the automorphic section
descends to an integer-weight Siegel or paramodular form on \(\Gamma\).

\paragraph{CHL ladder.}
On the \(K3\times E\) CHL ladder \(N\in\{1,2,3,4,6\}\), the canonical
constant-term table is
\[
  (c_N(0))=(10,8,6,4,2),
  \qquad
  \bigl(\kBKM(\Phi_N)\bigr)
  =c_N(0)/2=(5,4,3,2,1).
\]
The scalar-square denominator has doubled weight
\((10,8,6,4,2)\).  The primitive BKM denominator has weight
\((5,4,3,2,1)\).  At \(N=1\),
\[
  \Phi_{10}^{\mathrm{un}}=\Delta_5^2,
  \qquad
  \kBKM(\Delta_5)=5.
\]
For \(N>1\), identifying a scalar physical denominator with the square
of a primitive BKM denominator is a separate normalisation theorem.

\subsection*{Metaplectic boundary}

\paragraph{Definition.}
\(\CYhalf\subset\CYMp\) is the full subcategory whose objects carry a
nontrivial multiplier on the metaplectic kernel.  The automorphic target
must be named by its arithmetic group, cover, and character.  In the
Gritsenko-Clery diagonal-divisor catalogue the half-integer entries are
the rows of weights \(1/2\) and \(3/2\), with multiplier orders \(8\)
and \(4\), respectively.  They live on their specified paramodular
covers and are not inputs to a general CY\(_d\)-to-chiral functor.

\paragraph{Geometric scope.}
A metaplectic automorphic form supplies a target line for a super-VOA
only after the CY\(_2\) category, Mukai lattice, arithmetic subgroup,
and multiplier have been fixed.  The assignment
\[
  \Phisup:\CYMp\longrightarrow\SVOA
\]
is therefore a partially constructed functor: on \(\CYint\) it restricts
to the bosonic CY\(_2\) output
\[
  \Phi_2:\CY_2\to E_2\text{-}\ChirAlg\subset\VOA\subset\SVOA,
\]
while on \(\CYhalf\) it requires the multiplier-compatible fermionic
sector and the square-root automorphic line.  For \(d\ge3\) the
CY-to-chiral output on a curve is \(E_1\); an \(E_2\) structure is
available only after passing to the centre or double under the stated
hypotheses.

\subsection*{Complete definition}

\paragraph{Theorem within the stated axioms.}
\(\CYMp\) is the \((\infty,1)\)-category whose objects are CY\(_2\)
categories equipped with a rank-four Mukai lattice, an arithmetic
monodromy subgroup of \(\Sp_4(\Z)\), a metaplectic lift, and a multiplier
character.  Its morphisms are CY\(_2\) Fourier-Mukai kernels preserving
the lifted period data.  The forgetful map to \(\CY_2\) is a
\(B(\Z/2)\)-gerbe classified by \(w_2^{\mathrm{per}}\).  The integer
slice \(\CYint\) carries descended integer-weight automorphic forms,
including the CHL ladder above.  The boundary slice \(\CYhalf\) carries
forms with nontrivial metaplectic multiplier; each such object must
record its cover and character explicitly.

\paragraph{Primary literature.}
Shimura 1975, \emph{Annals of Mathematics} 102; Howe 1979,
\emph{Proceedings of Symposia in Pure Mathematics} 33; Borcherds 1986,
\emph{Proceedings of the National Academy of Sciences} 83; Kac 1998,
\emph{Vertex Algebras for Beginners}; Giraud 1971,
\emph{Cohomologie non abelienne}; To\"en-Vaquie 2007,
\emph{Annales Scientifiques de l'Ecole Normale Superieure} 40;
Kontsevich-Soibelman 2009; Mukai 1988, \emph{Inventiones Mathematicae}
94; Cl\'ery-Gritsenko 2008, \emph{Acta Arithmetica} 133; Lorgat 2020,
arXiv:2004.01676.
